\section{\algo's Algorithm Design}
\label{sec:selection}
% In this section, we propose an automatic option selection mechanism to guide users in selecting suitable options for Benne. The mechanism consists of four functions that help users choose the most appropriate summarizing data structure, window model, outlier detection mechanism, and refinement strategy based on their primary performance objective and input data stream characteristics. The primary performance objective can be Accuracy, Efficiency, or Balance.
In this section, we discuss the completed algorithm design of \algo. Through flexible adaptability among varying design choices based on the varying stream workload characteristics and user-defined optimization targets, \algo can better utilize the advantages of every different design choices and thus guarantee a better clustering behavior under the specific optimization target.

\subsection{Workflow of \algo}
The \algo algorithm follows a two-phase execution strategy: online phase and offline phase. The adaptability of \algo comes from the configurability of its components, including the summarizing data structure ($struc.$), window model ($win.$), outlier detection mechanism ($out.$), and refinement strategy ($ref.$). According to the evolving stream characteristics detected from the collected stream samples $queue$ and with the predefined threshold values $thresholds$, This adaptability allows \algo to cater to the specific needs of various applications and data streams.
Algorithm~\ref{alg:algo} presents the high-level execution flow of \algo. 

During the online phase, the algorithm processes the input data stream by following these steps:
\emph{Window Function} (Line 2): The \texttt{Window Fun.} manages the data points in the summarizing data structure ($s$) according to the selected window model ($win.$). This function ensures that the clustering process focuses on the most relevant data points and adapts to the evolving data distributions.
\emph{Outlier Detection} (Lines 3-7): If the outlier detection mechanism ($out.$) is enabled, the \texttt{Outlier Fun.} is invoked to determine whether the current input point ($p$) is an outlier. If the point is not an outlier, it is inserted into the summarizing data structure ($s$) and updated it accordingly (Line 6). If the outlier detection mechanism is not enabled, the input point is directly inserted into the summarizing data structure (Line 8).
\emph{Incremental Refinement} (Lines 9-10): If the refinement strategy ($ref.$) is set to \incrementalRefinement, the algorithm performs the \texttt{Refine Fun.} as new data points are processed. This approach ensures that the clustering model remains up-to-date and adaptive to changes in the data distribution but may require more frequent updates, increasing computational overhead.

After processing the entire input data stream, the algorithm enters the offline phase:
\emph{One-shot Refinement} (Lines 11-12): If the refinement strategy ($ref.$) is set to \oneshotRefinement, the algorithm performs the \texttt{Refine Fun.} in a single pass at the end of the stream processing. This approach reduces the computational overhead by performing refinement less frequently but may result in a less adaptive model if changes in the data distribution are not captured in a timely manner.
By setting different types of $struc.$, $win.$, $out.$, and $ref.$, the \algo algorithm can be easily reconfigured into various variants, allowing it to adapt to different scenarios and data stream requirements.

% \begin{algorithm}[h]
% \scriptsize
% \caption{Execution flow of \algo.}
% \label{alg:algo}
%     \KwData{$p$ \tcp{Input point}}
%     \KwData{$s$ \tcp{Summarizing data structure}}
%     \SetKwInput{KwData}{Input}
%     \KwData{$struc.$ \tcp{Selected type of summarizing data structure}}   
%     \KwData{$win.$ \tcp{Selected type of window model}}
%     \KwData{$out.$ \tcp{Selected type of outlier detection mechanism}}
%     \KwData{$ref.$ \tcp{Selected type of refinement strategy}}
%     \tcp{Online Phase}
%     \While{!stop processing of input streams}{
%         \texttt{Window Fun. (...)}\;
%         \eIf{\text{$out.$ != \nod}}{
%             $b$ $\gets$ \texttt{Outlier Fun. (...)}\;
%             \If{\text{$b$ = false}}{
%                 \texttt{Insert Fun. (...)\tcp{Insert $p$ to $s$ and update $s$}}            
%             }
%         }{
%             \texttt{Insert Fun. (...)\tcp{Insert $p$ to $s$ and update $s$}}
%         }
%         \If{\text{$ref.$ = \incrementalRefinement}}{
%             \texttt{Refine Fun. ($ref.$)};
%         }
%     }
%     \tcp{Offline Phase}
%     \If{\text{$ref.$ = \oneshotRefinement}}{
%         \texttt{Refine Fun. ($ref.$)};
%     }
% \end{algorithm}

\subsection{Flexible Design Choice Configuration}
\begin{algorithm}[h]
\scriptsize
\caption{Execution flow of \algo.}
\label{alg:algo}
    \KwData{$p$ \tcp{Input point}}
    \KwData{$c$ \tcp{Temporal Clusters}}
    \KwData{$o$ \tcp{User's primary performance objective}}
    \KwData{$threshold$ \tcp{Identifying characteristics' changes}}
    \KwData{$queue$ \tcp{Batch of stream for detection}}
    \KwData{$struc.$ \tcp{Type of summarizing data structure}}   
    \KwData{$win.$ \tcp{Type of window model}}
    \KwData{$out.$ \tcp{Type of outlier detection mechanism}}
    \KwData{$ref.$ \tcp{Type of refinement strategy}}
    \tcp{Online Phase}
    Initializing parameters and design choices \;
    \While{!stop processing of input streams}{
      \eIf{\text{$queue$ is full}}{
            $struc._{old}$ = $struc.$ \;
            $characteristics$ = \texttt{Auto-detection Fun.(...)} \;
            $struc.$, $win.$, $out.$, $ref.$ = \texttt{Auto-selection Fun.(...)} \;
            \If{\text{$struc.$ != $struc._{old}$}}{
                Extract temporal clusters $c$ from old Summarizing Data Structure $s'$\;
                Initialize new Summarizing Data Structure $s$ based on $struc.$ and $c$\;
            }
            Empty $queue$ \;
        }{
            $queue$.push($p$) \;
        } 
        \texttt{Window Fun. (...)}\;
        \eIf{\text{$out.$ != \nod}}{
            $b$ $\gets$ \texttt{Outlier Fun. (...)}\;
            \If{\text{$b$ = false}}{
                \texttt{Insert Fun. (...)\tcp{Insert $p$ to $s$ and update $s$}}            
            }
        }{
            \texttt{Insert Fun. (...)\tcp{Insert $p$ to $s$ and update $s$}}
        }
        \If{\text{$ref.$ = \incrementalRefinement}}{
            \texttt{Refine Fun. ($ref.$)};
        }
    }
    \tcp{Offline Phase}
    \If{\text{$ref.$ = \oneshotRefinement}}{
        \texttt{Refine Fun. ($ref.$)};
    }
\end{algorithm}

The \emph{Window Fun.} (Algorithm~\ref{alg:algo_window}) applies different computing logic depending on the selected type of window models. The function maintains a counter $c$ to keep track of the number of processed data points and performs different actions based on the selected window model. 
\emph{Landmark Window} (Lines 4-8): If the \landmark window model is selected, the function checks whether the counter $c$ has exceeded the current landmark $m$. If so, the clustering results from the summarizing data structure $s$ are sunk (i.e., stored or output), and all clustering information is cleared. Then, the function updates the current landmark $m_{next}$ for the next window. This model is suitable for scenarios where data distributions change significantly over time, and the algorithm needs to adapt to these changes by periodically resetting the clustering information.
\emph{Damped Window} (Lines 9-10): In the \damped window model, the function updates the weights of data points and clusters in the summarizing data structure $s$ using the decay parameters $dc = (\alpha, \lambda)$. This model is designed for cases where older data points should have less influence on the clustering results. By applying a decay function, the algorithm can gradually forget outdated information and adapt to the current data distribution.
\emph{Sliding Window} (Lines 11-13): If the \sliding window model is chosen, the function checks whether the counter $c$ has exceeded the sliding window size $ws$. If so, the earliest data point in the window is removed from the summarizing data structure $s$. This model maintains a fixed-size window of the most recent data points, ensuring that the algorithm focuses on the current data distribution and disregards older data points that may no longer be relevant.

\begin{algorithm}[h]
\scriptsize
\caption{Window Fun. of \algo.}
\label{alg:algo_window}
    \SetKwFunction{FMain}{Window Fun.}
    \SetKwProg{Fn}{Function}{:}{}
    \tcc{$c$: counter (initialized with 0), $m$: landmark, $ws$: sliding window size, $dc = (\alpha, \lambda)$: decay parameters}
    \Fn{\FMain{$s$, $m$, $ws$, $dc$}}{
        $c$++\;
        \eIf{\text{$win.$ = \landmark }}{
            \If{\text{$c >$ $m$}}{
                Sink the clustering results from $s$\;
                Clear all of the clustering information\;
                $m$ $\gets$ $m_{next}$; \tcc{update current landmark}
            }                  
        }{\eIf{\text{$win.$ = \damped}}{
            Update weight with $dc$\;
            }(\tcp*[h]{\text{$win.$ = \sliding}}){
                \If{\text{$c >$  $ws$}}{
                    Remove the earliest point from window\;
                }                
            }        
               
        } 
    }    
\end{algorithm}

If outlier detection is enabled, the \emph{Outlier Fun.} (Algorithm~\ref{alg:algo_outlier}) is invoked. \emph{Outlier Fun.}  plays a crucial role in the outlier detection process of the \algo algorithm. It is responsible for detecting and managing outliers using different mechanisms based on the selected type of outlier detection ($out.$). The following is a more detailed explanation of the steps involved in the \emph{Outlier Fun.}:
\emph{Buffer Optimization} (Lines 2-8): If the selected outlier detection mechanism includes buffer or buffer with timer (\buffer or \buffertimer), the function checks if the input point ($p$) is an outlier. If $p$ is an outlier, it is inserted into the outlier buffer ($Buffer$) and assigned to the closest cluster ($cl$). The buffer cluster is then updated accordingly. After updating the cluster, the function checks if the cluster is dense enough (i.e., the density exceeds the specified threshold $d$). If the cluster is dense enough, it is moved from the outlier buffer to the summarizing data structure ($s$).
\emph{Regular Check} (Lines 9-24): The regular check is triggered at predefined intervals to assess the clusters in the summarizing data structure and the outlier buffer. For each cluster in the summarizing data structure, the function checks if the cluster is dense enough (density threshold $d$) and, if a timer-based mechanism is used (\timer or \buffertimer), if the cluster is active enough (timer threshold $t$). If a cluster is neither dense nor active, it is either moved to the outlier buffer (if enabled) or removed from the summarizing data structure.
\emph{Buffer Timer Check} (Lines 20-24): If the selected outlier detection mechanism is \buffertimer, the function checks the activity of each cluster in the outlier buffer based on the timer threshold ($t$). If a cluster in the buffer is not active, it is removed from the buffer.
\emph{Outlier Determination} (Line 25): The function returns whether the input point ($p$) is an outlier or not, based on the results of the previous steps.
By employing different outlier detection mechanisms, \algo is capable of efficiently handling various types of data streams and application requirements.

\begin{algorithm}[t]
\scriptsize
\caption{Outlier Fun. of \algo.}
\label{alg:algo_outlier}
    \SetKwFunction{FMain}{Outlier Fun.}
    \SetKwProg{Fn}{Function}{:}{}
    \tcc{$Buffer$: outlier buffers, $t$: timer threshold, $d$: density threshold}
    \Fn{\FMain{$p$, $s$, $Buffer$, $t$, $d$}}{ 
        \If{$out.$ = (\buffer \textbf{ Or } \buffertimer)}{
            \tcc{use buffer optimization}
            \If{($p$ Is Outlier) = true}{
                Select cluster $cl$ closest to $p$ from $Buffer$\;
                Insert $p$ to $cl$ and update $cl$\;
                bool $dense \gets$ Check ($cl$)\;
                \tcc{check if $cl$ is dense enough}
                \If{$dense$ = true}{
                    Move $cl$ from $Buffer$ to $s$\;
                }
            }
        }        
        \If{REGULAR CHECK TRIGGERED}{
            $cls \gets$ Extract($s$);\tcc{extract all clusters from $s$}
            \For{$cl \in cls$}{
                bool $active$ = false\;
                bool $dense \gets$ Check ($cl, d$)\;
                \tcc{check if $cl$ is dense enough}
                \If{$out.$ = (\timer \textbf{ Or } \buffertimer)}{
                    \tcc{use timer optimization}
                    bool $active \gets$ Check($cl.timer, t$)\;
                    \tcc{check if $cl$ is active enough }
                }
                \If{$dense$ = false \textbf{ and } $active$ = false}{
                    \If{BUFFER ENABLED}{
                        Insert $cl$ to $Buffer$\;
                    }
                    Remove $cl$ from $s$\; 
                }
            }   
            \If{\buffertimer}{
                \For{$cl \in Buffer$}{
                     bool $active \gets$ Check($cl.timer, t$)\;
                     \If{$active$ = false}{
                        Remove $cl$ from $Buffer$\; 
                    }
                }
            } 
        }
        Return ($p$ Is Outlier)\;
    }  
\end{algorithm}

The input point will then be added to the summarizing data structure if it is not identified as an outlier or if the algorithm does not use any outlier detection mechanisms.
In the \emph{Insert Fun.} (Algorithm~\ref{alg:algo_insert}), the input point $p$ is inserted into the summarizing data structure $struct.$. The function can adapt to different types of summarizing data structures, specifically those belonging to the hierarchical category and the partitional category. This adaptability of the \emph{Insert Fun.} allows \algo to effectively handle different types of summarizing data structures, catering to various needs and requirements. When the selected data structure belongs to the hierarchical category, the function inserts the input point $p$ into the appropriate cluster within the hierarchical data structure and updates the hierarchy accordingly (Lines 2-3). On the other hand, if the selected data structure belongs to the partitional category, the function first identifies the closest cluster for the input point $p$ and then inserts it into the identified cluster (Lines 4-5). After the insertion, the summarizing data structure is updated accordingly. 

\begin{algorithm}[h]
\scriptsize
\caption{Insert Fun. of \algo.}
\label{alg:algo_insert}
\SetKwFunction{FMain}{Insert Fun.}
\SetKwProg{Fn}{Function}{:}{}
\tcc{$p$: input, $struc.$: summarizing data structure}
\Fn{\FMain{$p$, $struc.$}}{
    \eIf{\text{$struc.$ belongs to hierarchical category}}{
    \tcc{Insert $p$ into the hierarchical data structure and update it}
        Insert $p$ into the appropriate cluster in the hierarchy;
        Update the hierarchy based on the insertion;
    }{
    \tcc{Insert $p$ into the partitional data structure and update it}
        Find the closest cluster for $p$;
        Insert $p$ into the closest cluster;
        Update the partitional data structure based on the insertion;
    }
}
\end{algorithm}

The \emph{Refine Fun.} (Algorithm~\ref{alg:algo_refine}) is responsible for refining the clustering results based on the selected refinement strategy ($ref.$). \algo provides two distinct refinement strategies: \incrementalRefinement and \oneshotRefinement, catering to different application requirements and data stream characteristics.
When \incrementalRefinement is applied, \algo continuously extracts the current clustering model from the summarizing data structure $s$. Then, it applies a suitable batch clustering algorithm to update the clustering model. After that, the function updates $s$ with the new clustering model and removes any outdated or redundant data points or clusters from $s$. This process ensures that the model remains up-to-date and adapts to the evolving data distribution.
When \oneshotRefinement is applied, it extracts all clusters and data points from the summarizing data structure $s$. It then applies a batch clustering algorithm to refine the clustering results. Finally, the function updates $s$ with the refined clustering results, providing a more accurate representation of the data distribution at the end of the processing. 

\begin{algorithm}[h]
\scriptsize
\caption{Refine Fun. of \algo.}
\label{alg:algo_refine}
\SetKwFunction{FMain}{Refine Fun.}
\SetKwProg{Fn}{Function}{:}{}
\tcc{$ref.$: Selected type of refinement strategy, $s$: Summarizing data structure}
\Fn{\FMain{$ref.$, $s$}}{
    \eIf{\text{$ref.$ = \incrementalRefinement}}{
        \tcc{Apply incremental refinement}
        Extract the current clustering model from $s$;
        Apply batch clustering algorithm to update the clustering model;
        Update $s$ with the new clustering model;
        Remove outdated or redundant data points or clusters from $s$;
    }{
        \tcc{Apply one-shot refinement}
        Extract all clusters and data points from $s$;
        Apply batch clustering algorithm to refine the clustering results;
        Update $s$ with the refined clustering results;
    }
}
\end{algorithm}



\subsection{Automatic Stream Characteristics Detection}
\begin{algorithm}[h]
\scriptsize
\caption{Auto-detection Func. of \algo}
\label{alg:auto_detection}
\SetKwFunction{FMain}{Auto-detection Func.}
\SetKwProg{Fn}{Function}{:}{}
\tcc{$batch$: a small batch of input stream data stored for detecting stream characteristic changes, $centers$: previous clustering centers of the stream data, $thresholds$: threshold values for identifying the characteristic}
\Fn{\FMain{$batch, centers, threshold$}}{
    \tcc{Initialize stream characteristics}
    \text{$characteristics$ = null}\;
    \text{$highDimData$ = $var$ = $numOfOutliers$= 0}\;
    \text{$samples$ = UniformSampling($batch$)}\;
    \text{$new\_center$ = mean($samples$)}\;
    \For{$s \in samples$}{
        \If{$s.dim  > thresholds.dim$}{
            $highDimData$ += 1 \;
        }
        $var$ = UpdateVariance($var, s, new\_centers$)\;
        $centers$ += $new\_center$ \;
        \If{$min(dist(s, centers))  > thresholds.dist$}{
            $numOfOutliers$ += 1 \;
        }
    }
    \eIf{$highDimData > num(samples) / 2$}{
        $characteristics.HighDimension = true$ \;
    }{
        $characteristics.HighDimension = false$ \;
    }
    \eIf{$var > thresholds.variance$}{
        $characteristics.HighConceptDrift = true$ \;
    }{
        $characteristics.HighConceptDrift = false$ \;
    }
    \eIf{$numOfOutliers > thresholds.outlier$}{
        $characteristics.ManyOutliers = true$ \;
    }{
        $characteristics.ManyOutliers = false$ \;
    }
    Return $characteristics$; 
}
\end{algorithm}

\subsection{Automatic Module Selection}
\begin{algorithm}[h]
\scriptsize
\caption{Auto-selection Func. of \algo}
\label{alg:auto_selection}
\SetKwFunction{FMain}{Auto-Selection Func.}
\SetKwProg{Fn}{Function}{:}{}
\tcc{$objective$: User's primary performance objective (Accuracy, Efficiency, or Balance); $data$: Input data stream; $characteristics$: Stream characteristics detected by auto-detect Func.}
\Fn{\FMain{$objective$, $data$, $characteristics$}}{
    \tcc{Initialize  four modular design components}
    \text{$struc.$= $win.$ = $out.$= $ref.$ = null}\;
    \eIf{$objective$ = Accuracy}{
        \tcc{Select modules with high accuracy based on input data stream characteristics}
        \eIf{$characteristics.HighConceptDrift = true$}{
            $struc.$ = \MCs \;
        }{
            $struc.$ = \MS \;
        }
        \eIf{$characteristics.ManyOutliers = true$}{
            $win.$ = \landmark \;
            $out.$ = \buffertimer \;
        }{
            $win.$ = \damped \;
            \eIf{$characteristics.HighDimension = true$}{
                $out.$ = \buffer \;
            }{
                $out.$ = \buffertimer \;
            }
        }
        $ref.$ = \incrementalRefinement \;
    }
    {
        \eIf{$objective$ = Efficiency}{
            \tcc{Choose a structure with high efficiency based on input data stream characteristics}
            \eIf{$characteristics.HighConceptDrift = true$}{
                $struc.$ = \DPT \;
                $win.$ = \landmark \;
            }{
                $struc.$ = \Grids \;
                $win.$ = \sliding \;
            }
            \eIf{$characteristics.ManyOutliers = true$}{
                $out.$ = \od \;
            }{
                $out.$ = \nod \;
            } 
            $ref.$ = \norefinement \;
        }{
            \tcc{Choose a structure with balanced accuracy and efficiency based on input data stream characteristics}
            \eIf{$characteristics.HighConceptDrift = true$}{
                $struc.$ = \CFT \;
            }{
                $struc.$ = \CoreT \;
            }
            $win.$ = \landmark \;
            $out.$ = \timer \;
            $ref.$ = \oneshotRefinement \;
        }
    }
    Return $struc.$, $win.$, $out.$, $ref.$; 
}
\end{algorithm}
% \subsection{Selecting Summarizing Data Structure}
% The SelectStructure function (Algorithm~\ref{alg:select_structure}) chooses the most suitable summarizing data structure based on the user's primary performance objective and the characteristics of the input data stream. Based on our observations, different summarizing data structures have different trade-offs between accuracy and efficiency. The function analyzes the input data stream to select the best structure that matches the user's preferences.

% If the user prefers accuracy, it selects a structure with high accuracy like \CoreT or \DPT. However, the function also considers the input data stream's characteristics to make a more informed decision. For example, if the input data has a high rate of change, \DPT may be more suitable due to its flexible updating mechanism.

% If the user prefers efficiency, it selects a structure with high efficiency like \CFT or \Grids. The function takes into account the input data stream's characteristics to select the most efficient structure for the given data. For example, if the input data has a low dimensionality and uniform distribution, \Grids might be a better choice due to its efficient distance calculation elimination.

% If the user seeks a balance between accuracy and efficiency, it selects a structure with balanced performance like \DPT or \CFT. The function considers the input data stream's characteristics to select the most suitable structure that provides a good balance between accuracy and efficiency.

% \begin{algorithm}[h]
% \scriptsize
% \caption{SelectStructure Function}
% \label{alg:select_structure}
% \SetKwFunction{FMain}{SelectStructure}
% \SetKwProg{Fn}{Function}{:}{}
% \tcc{$objective$: User's primary performance objective (Accuracy, Efficiency, or Balance); $data$: Input data stream characteristics}
% \Fn{\FMain{$objective$, $data$}}{
% \eIf{$objective$ = Accuracy}{
%     \tcc{Choose a structure with high accuracy based on input data stream characteristics}
%     \eIf{High rate of change in $data$}{
%         Return \DPT
%     }{
%         Return \CoreT
%     }
%     }{
%     \eIf{$objective$ = Efficiency}{
%     \tcc{Choose a structure with high efficiency based on input data stream characteristics}
%     \eIf{Low dimensionality and uniform distribution in $data$}{
%         Return \Grids
%     }{
%         Return \CFT
%     }
%     }{
%     \tcc{Choose a structure with balanced accuracy and efficiency based on input data stream characteristics}
%     \eIf{High rate of change in $data$}{
%         Return \DPT
%     }{
%         Return \CFT
%         }
%     }
%     }
% }
% \end{algorithm}

% \subsection{Selecting Window Model}

% The SelectWindowModel function (Algorithm~\ref{alg:select_window_model}) selects the most appropriate window model based on the user's primary performance objective and input data stream characteristics. Each window model has its strengths and weaknesses concerning accuracy and efficiency. The function takes into account both the user's preferences and the input data stream characteristics to select the most suitable window model.

% If the user prefers accuracy, it chooses a model with high accuracy like \landmark or \damped. The function also considers the input data stream's characteristics to make a more informed decision. For example, if the input data has a high rate of change, \damped may be more suitable due to its adaptive forgetting factor.

% If the user prefers efficiency, it chooses a model with high efficiency like \sliding. The function takes into account the input data stream's characteristics to select the most efficient window model for the given data. For example, if the input data has a relatively stable distribution, \sliding might be a better choice due to its fixed window size, which reduces the number of clusters for maintenance.

% If the user seeks a balance between accuracy and efficiency, it analyzes the input data stream and selects a suitable window model based on its characteristics and the desired balance. The function might choose \landmark or \damped for a balance favoring accuracy, or \sliding for a balance favoring efficiency. The specific choice would depend on the input data stream's characteristics and the desired balance of performance.

% \begin{algorithm}[h]
% \scriptsize
% \caption{SelectWindowModel Function}
% \label{alg:select_window_model}
% \SetKwFunction{FMain}{SelectWindowModel}
% \SetKwProg{Fn}{Function}{:}{}
% \tcc{$objective$: User's primary performance objective (Accuracy, Efficiency, or Balance); $stream$: Input data stream characteristics}
% \Fn{\FMain{$objective$, $stream$}}{
%     \eIf{$objective$ = Accuracy}{
%     \tcc{Choose a model with high accuracy based on input data stream characteristics}
%     \eIf{High rate of change in $stream$}{
%         Return \damped
%     }{
%         Return \landmark
%     }
%     }{
%     \eIf{$objective$ = Efficiency}{
%         \tcc{Choose a model with high efficiency based on input data stream characteristics}
%     \eIf{Stable distribution in $stream$}{
%         Return \sliding
%     }{
%         Analyze $stream$ and choose a suitable window model
%     }
%     }{
%         \tcc{Choose a suitable window model based on stream characteristics and balance}
%         Analyze $stream$ and select a window model that balances accuracy and efficiency
%     }
%     }
% }
% \end{algorithm}

% \subsection{Selecting Outlier Detection Mechanism}

% The SelectOutlierDetection function (Algorithm~\ref{alg:select_outlier_detection}) picks the best outlier detection mechanism based on the user's primary performance objective and input data stream characteristics. Each outlier detection mechanism has its strengths and weaknesses concerning accuracy and efficiency. The function takes into account both the user's preferences and the input data stream characteristics to select the most suitable outlier detection mechanism.

% If the user prefers accuracy, it chooses a mechanism with high accuracy like \buffer or \buffertimer. The function considers the input data stream's characteristics to make a more informed decision. For example, if the input data has a high rate of change, \buffertimer may be more suitable due to its combination of buffer and timer mechanisms, which can handle the evolving data distribution more effectively.

% If the user prefers efficiency, it chooses a mechanism with high efficiency like \timer. The function takes into account the input data stream's characteristics to select the most efficient outlier detection mechanism for the given data. For example, if the input data has a relatively stable distribution, \timer might be a better choice due to its lower maintenance overhead.

% If the user seeks a balance between accuracy and efficiency, it analyzes the input data stream and selects a suitable outlier detection mechanism based on its characteristics and the desired balance. The function might choose \buffer, \timer, or \buffertimer, depending on the input data stream's characteristics and the desired balance of performance. The specific choice would depend on the evolving nature of the input data stream and the user's tolerance for outlier detection errors.

% \begin{algorithm}[h]
% \scriptsize
% \caption{SelectOutlierDetection Function}
% \label{alg:select_outlier_detection}
% \SetKwFunction{FMain}{SelectOutlierDetection}
% \SetKwProg{Fn}{Function}{:}{}
% \tcc{$objective$: User's primary performance objective (Accuracy, Efficiency, or Balance); $stream$: Input data stream characteristics}
% \Fn{\FMain{$objective$, $stream$}}{
% \eIf{$objective$ = Accuracy}{
% \tcc{Choose a mechanism with high accuracy based on input data stream characteristics}
% \eIf{High rate of change in $stream$}{
% Return \buffertimer
% }{
% Return \buffer
% }
% }{
% \eIf{$objective$ = Efficiency}{
% \tcc{Choose a mechanism with high efficiency based on input data stream characteristics}
% \eIf{Stable distribution in $stream$}{
% Return \timer
% }{
% Analyze $stream$ and choose a suitable outlier detection mechanism
% }
% }{
% \tcc{Choose a suitable outlier detection mechanism based on stream characteristics and balance}
% Analyze $stream$ and select an outlier detection mechanism that balances accuracy and efficiency
% }
% }
% }
% \end{algorithm}

% \subsection{Selecting Refinement Strategy}

% The SelectRefinementStrategy function (Algorithm~\ref{alg:select_refinement_strategy}) selects the most suitable refinement strategy based on the user's primary performance objective and input data stream characteristics. If the user prefers accuracy, it chooses a strategy with high accuracy like \incrementalRefinement or \oneshotRefinement. If the user prefers efficiency, it may choose not to apply refinement. If the user seeks a balance between accuracy and efficiency, it analyzes the input data stream and selects a suitable refinement strategy based on its characteristics and the desired balance.

% The SelectRefinementStrategy function chooses the most suitable refinement strategy based on the user's primary performance objective and input data stream characteristics. The refinement strategies, such as \incrementalRefinement and \oneshotRefinement, can improve clustering accuracy but might affect the efficiency of the overall algorithm.

% If the user prefers accuracy, it chooses a strategy with high accuracy like \incrementalRefinement or \oneshotRefinement. The function considers the input data stream's characteristics to make a more informed decision. For example, if the input data has dynamic cluster formation, \incrementalRefinement may be more suitable due to its ability to adjust the clustering results iteratively.

% If the user prefers efficiency, it may choose not to apply any refinement strategy, as refinement can increase the computational cost and affect the overall efficiency of the algorithm.

% If the user seeks a balance between accuracy and efficiency, it analyzes the input data stream and selects a suitable refinement strategy based on its characteristics and the desired balance. The specific choice would depend on the input data stream's characteristics, the clustering algorithm's performance, and the user's tolerance for clustering errors. The function might choose \incrementalRefinement, \oneshotRefinement, or NoRefinement, depending on the input data stream's characteristics and the desired balance of performance.

% \begin{algorithm}[h]
% \scriptsize
% \caption{SelectRefinementStrategy Function}
% \label{alg:select_refinement_strategy}
% \SetKwFunction{FMain}{SelectRefinementStrategy}
% \SetKwProg{Fn}{Function}{:}{}
% \tcc{$objective$: User's primary performance objective (Accuracy, Efficiency, or Balance)}
% \tcc{$stream$: Input data stream characteristics}
% \Fn{\FMain{$objective$, $stream$}}{
%     \eIf{$objective$ = Accuracy}{
%         \tcc{Choose a strategy with high accuracy based on input data stream characteristics}
%         \eIf{Dynamic cluster formation in $stream$}{
%         Return \incrementalRefinement
%         }{
%         Return \oneshotRefinement
%         }
%     }{
%     \eIf{$objective$ = Efficiency}{
%         \tcc{Choose not to apply refinement for better efficiency}
%     Return NoRefinement
%     }{
%         \tcc{Choose a suitable refinement strategy based on stream characteristics and balance}
%         Analyze $stream$ and select a refinement strategy that balances accuracy and efficiency
%     }
%     }
% }
% \end{algorithm}

% \subsection{Decision Model}

% \textbf{Summarizing Data Structures.}
% In general, hierarchical summarizing data structure improves the clustering efficiency while the partitional structure is expected to bring better accuracy. The unique design of each specific structure, however, leads to further differences in their clustering behaviour. For example, although \CFT, \CoreT, \DPT, and \MCs promise high clustering accuracy due to the comprehensive information stored in their efficiency behaviour is quite different. Only \CFT and \DPT can guarantee high efficiency due to their flexible updating. \CoreT and \MS are slow due to frequently rebuilding the whole structure while \MCs is slow since it needs to scan every temporal cluster for data insertion. For \Grids, although it may not be as accurate as the other structures, it largely saves its computation time due to eliminating the distance calculation operation. 

% \textbf{Window Models.}
% Although these window models all try to prioritize the most recent data from streams, their specific strategies are quite different leading to different clustering behaviours. For a better clustering efficiency, \sliding may be the best choice since its fixed window size largely reduces the number of clusters for maintenance. On the contrary, although both \landmark and \damped may not be as efficient as the \sliding, they can bring in more accurate clustering results. Finally, for all three types of window models, their clustering behaviour is sensitive to their specific window configurations, which have not been studied in depth in the literature.

% \textbf{Outlier Detection Mechanisms.}
% The buffer and timer optimizations can be applied either separately or simultaneously, resulting in different decisions regarding the outlier detection mechanism. 
% The use of a buffer can greatly improve the clustering quality due to its better ability for handling \oef. However, maintaining the outlier buffer in real-time will also bring down the efficiency. 
% On the contrary, the use of the timer may be beneficial for both better clustering accuracy and efficiency, especially under \oef. 
% Using a buffer and timer together in outlier detection further complicates the clustering behaviour and motivates an in-depth study.

% \textbf{Refinement Strategies.}
% Refinement is optional and can be easily applied to refine online clustering results before outputting the result. The common wisdom is that it may further improve clustering quality. Unfortunately, there is no prior study about the real impact of this design aspect on both accuracy and efficiency. Even in recent years, some choose~\cite{StreamKM++:12} to involve this design aspect in their proposed \dsc algorithms, and some choose not~\cite{EDMStream:17}. We will empirically show its impact via extensive evaluations.